\documentclass[11pt]{ctexart}

\usepackage[a4paper,margin=1in]{geometry}
\usepackage{amsmath,amssymb,amsthm,bm,mathtools}
\usepackage{booktabs}
\usepackage{enumitem}
\usepackage{natbib}
\usepackage{xcolor}
\usepackage{listings}
\usepackage[colorlinks=true,linkcolor=blue,citecolor=blue,urlcolor=blue]{hyperref}

\lstset{
  language=R,
  basicstyle=\small\ttfamily,
  commentstyle=\color{gray},
  showstringspaces=false
}

\newtheorem{assumption}{Assumption}
\newtheorem{definition}{Definition}
\newtheorem{theorem}{Theorem}
\newtheorem{remark}{Remark}

\newenvironment{problemblock}[1]{%
  \subsection*{#1}
  \addcontentsline{toc}{subsection}{#1}
}{\medskip}

\newenvironment{prep}{%
  \paragraph{预备知识与可引用内容.}
  \begin{enumerate}[leftmargin=2em,itemsep=0.35em]
}{%
  \end{enumerate}
}

\newenvironment{solutionbox}{%
  \paragraph{我的解答.}
  \begin{quote}
  \textit{这里写 solution/proof。}
  \vspace{2.2cm}
  \end{quote}
}{}

\newcommand{\ri}[1]{\texttt{#1}}
\newcommand{\pr}{\operatorname{pr}}
\newcommand{\var}{\operatorname{var}}
\newcommand{\cov}{\operatorname{cov}}
\newcommand{\E}{\mathbb{E}}
\newcommand{\ind}{\perp\!\!\!\perp}
\newcommand{\sumn}{\sum_{i=1}^n}
\newcommand{\tran}{^{\mathsf T}}
\newcommand{\at}{\mathrm{a}}
\newcommand{\nt}{\mathrm{n}}
\newcommand{\cp}{\mathrm{c}}
\newcommand{\df}{\mathrm{d}}
\newcommand{\CACE}{\textsc{cace}}
\newcommand{\LATE}{\textsc{late}}
\newcommand{\TSLS}{\textsc{tsls}}
\newcommand{\ILS}{\textsc{ils}}
\newcommand{\RR}{\textsc{rr}}

\title{Homework 5 写作预备模板}
\author{姓名：\underline{\hspace{3cm}}}
\date{Due: May 27, 4:30 pm}

\begin{document}
\maketitle
\tableofcontents

\section*{使用说明}
这个文件只整理题面、出处、引用和预备知识；证明和实证分析部分留空，由我自己完成。
主要来源为 Peng Ding 的 \emph{A First Course in Causal Inference} 第 21--23 章
\citep{imbens1994identification,angrist1996identification}，以及作业 PDF 的 Problem Set 5。

\section{统一符号与基本假设}

\begin{definition}[Compliance type]
令 \(Z_i\in\{0,1\}\) 为 assignment 或 IV，\(D_i\in\{0,1\}\) 为实际接受的处理，
\(Y_i\) 为 outcome。关于 assignment 的潜在接受处理为 \(D_i(1),D_i(0)\)。
由 \(U_i=\{D_i(1),D_i(0)\}\) 定义四类人：
\[
U_i=
\begin{cases}
\at, & D_i(1)=1,\ D_i(0)=1 \quad \text{always taker},\\
\cp, & D_i(1)=1,\ D_i(0)=0 \quad \text{complier},\\
\df, & D_i(1)=0,\ D_i(0)=1 \quad \text{defier},\\
\nt, & D_i(1)=0,\ D_i(0)=0 \quad \text{never taker}.
\end{cases}
\]
\end{definition}

\begin{assumption}[Randomization]
\[
Z\ind \{D(1),D(0),Y(1),Y(0)\}.
\]
在 observational 版本中常替换为 conditional ignorability：
\[
Z\ind \{D(1),D(0),Y(1),Y(0)\}\mid X.
\]
\end{assumption}

\begin{assumption}[Monotonicity]
\[
D_i(1)\ge D_i(0), \quad \text{equivalently } \pr(U_i=\df)=0.
\]
\end{assumption}

\begin{assumption}[Exclusion restriction]
在 assignment 潜在结果记号下，always takers 和 never takers 满足 \(Y_i(1)=Y_i(0)\)。
在 double-index 记号下，排除限制写为
\[
Y_i(z,d)=Y_i(d),\quad z=0,1.
\]
\end{assumption}

\begin{definition}[CACE/LATE]
\[
\tau_{\cp}
=\E\{Y(1)-Y(0)\mid U=\cp\}
=\E\{Y(1)-Y(0)\mid D(1)=1,D(0)=0\}.
\]
在 double-index 记号下也写为
\[
\tau_{\cp}=\E\{Y(d=1)-Y(d=0)\mid U=\cp\}.
\]
\end{definition}

\begin{theorem}[CACE identification, Ding Chapter 21]
在 randomization、monotonicity 和 exclusion restriction 下，
\[
\tau_{\cp}
=
\frac{\E(Y\mid Z=1)-\E(Y\mid Z=0)}
{\E(D\mid Z=1)-\E(D\mid Z=0)}.
\]
对应样本估计量为 Wald/IV estimator：
\[
\hat\tau_{\cp}
=
\frac{\hat\tau_Y}{\hat\tau_D}.
\]
\end{theorem}

\section{Peng Ding 书中习题}

\begin{problemblock}{Chapter 21.1: Variance of the Wald estimator}
\paragraph{题面.}
Show that \(\var(\hat{\tau}_{\cp})=\infty\).

\paragraph{出处.}
Peng Ding, Chapter 21, Homework problems, ``Variance of the Wald estimator''
\(\langle\)label: \ri{problem::infinite-var-iv}\(\rangle\).

\begin{prep}
  \item Wald estimator:
  \[
  \hat{\tau}_{\cp}=\frac{\hat{\tau}_Y}{\hat{\tau}_D},
  \quad
  \hat{\tau}_Y=\bar Y_1-\bar Y_0,\quad
  \hat{\tau}_D=\bar D_1-\bar D_0.
  \]
  \item 书中提醒：即使 \(\tau_D>0\)，也可能有正概率出现
  \(\hat{\tau}_D=0\)，所以 ratio estimator 的有限样本方差可以为无穷。
  \item 可引用 weak IV 讨论：当 \(\tau_D\) 接近 \(0\) 时，
  \(\hat{\tau}_{\cp}\) 有较差有限样本性质，Wald-type confidence interval 可能覆盖率差。
  \item 参考引用：Wald estimator \citep{wald1940fitting}；CACE/LATE 框架
  \citep{imbens1994identification,angrist1996identification}。
\end{prep}

\begin{solutionbox}
\end{solutionbox}
\end{problemblock}

\begin{problemblock}{Chapter 21.6: Binary IV and ordinal treatment received}
\paragraph{题面.}
\citet{angrist1995two} considered binary \(Z\), ordinal treatment
\(D\in\{0,1,\ldots,J\}\), and outcome \(Y\). Potential treatment received is
\(D(1),D(0)\), and outcome has potential values \(Y(z,d)\).
Extend the IV assumptions and prove the theorem below.

\paragraph{Assumption to use.}
\[
Z\ind \{D(z),Y(z,d): z=0,1;\ d=0,1,\ldots,J\},
\quad
D(1)\ge D(0),
\quad
Y(z,d)=Y(d).
\]

\paragraph{The theorem to prove.}
\[
\frac{\E(Y\mid Z=1)-\E(Y\mid Z=0)}
{\E(D\mid Z=1)-\E(D\mid Z=0)}
=
\sum_{j=1}^J w_j
\E\{Y(j)-Y(j-1)\mid D(1)\ge j>D(0)\},
\]
where
\[
w_j=
\frac{\pr\{D(1)\ge j>D(0)\}}
{\sum_{j'=1}^J \pr\{D(1)\ge j'>D(0)\}}.
\]

\paragraph{出处.}
Peng Ding, Chapter 21, Homework problems, ``Binary IV and ordinal treatment received''
\(\langle\)label: \ri{para::iv-ordinal-treatment}\(\rangle\).

\begin{prep}
  \item 当 \(J=1\) 时，此 theorem 退化为 standard CACE/LATE identification。
  \item 可先写
  \[
  D(z)=\sum_{j=0}^J j\,1\{D(z)=j\},
  \quad
  Y(D(z))=\sum_{j=0}^J Y(j)1\{D(z)=j\}.
  \]
  \item 书中提示使用 Abel's lemma / summation by parts：
  \[
  \sum_{j=0}^J f_j(g_{j+1}-g_j)
  =
  f_Jg_{J+1}-f_0g_0-\sum_{j=1}^J g_j(f_j-f_{j-1}).
  \]
  \item 引用：ordinal treatment IV result 来自 \citet{angrist1995two}。
\end{prep}

\begin{solutionbox}
\end{solutionbox}
\end{problemblock}

\begin{problemblock}{Chapter 21.7: Data analysis: a flu shot encouragement design}
\paragraph{题面.}
The dataset \(\ri{fludata.txt}\) is from a randomized encouragement design of
\citet{mcdonald1992effects}, also re-analyzed by \citet{hirano2000assessing}.
Analyze the data with and without adjusting for covariates.

\paragraph{数据变量.}
\[
\begin{array}{ll}
\texttt{assign} & \text{binary encouragement to receive the flu shot},\\
\texttt{receive} & \text{binary indicator for receiving the flu shot},\\
\texttt{outcome} & \text{binary outcome for flu-related hospitalization},\\
\texttt{age} & \text{age},\\
\texttt{sex} & \text{sex},\\
\texttt{race} & \text{race},\\
\texttt{copd} & \text{chronic obstructive pulmonary disease},\\
\texttt{dm} & \text{diabetes},\\
\texttt{heartd} & \text{heart disease},\\
\texttt{renal} & \text{renal disease},\\
\texttt{liverd} & \text{liver disease}.
\end{array}
\]

\paragraph{出处.}
Peng Ding, Chapter 21, Homework problems, ``Data analysis: a flu shot encouragement design''
\(\langle\)label: \ri{problem::flushot}\(\rangle\).

\begin{prep}
  \item Without covariates: use the Wald estimator
  \[
  \hat\tau_{\cp}
  =
  \frac{\bar Y_{Z=1}-\bar Y_{Z=0}}{\bar D_{Z=1}-\bar D_{Z=0}}.
  \]
  \item With covariates under CRE: use Lin-style adjustment separately for \(Y\) and \(D\),
  then take the ratio:
  \[
  \hat{\tau}_{\cp,L}
  =
  \frac{\hat{\tau}_{Y,L}}{\hat{\tau}_{D,L}}.
  \]
  \item Recommended reporting: first stage \(\hat\tau_D\), ITT \(\hat\tau_Y\),
  CACE estimate, standard error/confidence interval, and assumptions
  (randomization, monotonicity, exclusion restriction).
  \item Local data path:
  \[
  \ri{/Users/yitwah/homework/fludata.txt}.
  \]
\end{prep}

\begin{solutionbox}
\end{solutionbox}
\end{problemblock}

\begin{problemblock}{Chapter 22.2: Risk ratio for compliers}
\paragraph{题面.}
With a binary outcome, define the risk ratio for compliers:
\[
\RR_{\cp}
=
\frac{\pr\{Y(1)=1\mid U=\cp\}}
{\pr\{Y(0)=1\mid U=\cp\}}.
\]
Show that under Assumptions randomization--exclusion restriction,
\[
\RR_{\cp}
=
\frac{\E(DY\mid Z=1)-\E(DY\mid Z=0)}
{\E\{(D-1)Y\mid Z=1\}-\E\{(D-1)Y\mid Z=0\}}.
\]

\paragraph{出处.}
Peng Ding, Chapter 22, Homework problems, ``Risk ratio for compliers''
\(\langle\)label: \ri{hw::rr-compliers}\(\rangle\).

\begin{prep}
  \item Use Theorem 22.1, disentangling mixture means:
  \[
  \pi_{\nt}=\pr(D=0\mid Z=1),\quad
  \pi_{\at}=\pr(D=1\mid Z=0),\quad
  \pi_{\cp}=\E(D\mid Z=1)-\E(D\mid Z=0).
  \]
  \item The complier potential outcome means are identifiable:
  \[
  \mu_{1\cp}
  =
  \pi_{\cp}^{-1}\{\E(DY\mid Z=1)-\E(DY\mid Z=0)\},
  \]
  \[
  \mu_{0\cp}
  =
  \pi_{\cp}^{-1}
  [\E\{(1-D)Y\mid Z=0\}-\E\{(1-D)Y\mid Z=1\}].
  \]
  \item Since \((D-1)Y=-(1-D)Y\), the denominator in the target formula is
  \(\E\{(1-D)Y\mid Z=0\}-\E\{(1-D)Y\mid Z=1\}\).
  \item Citation: mixture identification follows the IV framework of
  \citet{imbens1994identification,angrist1996identification}.
\end{prep}

\begin{solutionbox}
\end{solutionbox}
\end{problemblock}

\begin{problemblock}{Chapter 22.7: Bounds on the average causal effect on the whole population}
\paragraph{题面.}
Using the double-index notation \(Y(d)\), define the average causal effect of the
treatment received on the whole population:
\[
\delta=\E\{Y(d=1)-Y(d=0)\}.
\]
Let
\[
m_{du}=\E\{Y(d)\mid U=u\},\quad u=\at,\nt,\cp.
\]
Then
\[
\delta=\sum_{u=\at,\nt,\cp}\pi_u(m_{1u}-m_{0u}).
\]
Section 22 identifies
\(\pi_{\at},\pi_{\nt},\pi_{\cp},m_{1\at},m_{0\nt},m_{1\cp},m_{0\cp}\),
but \(m_{0\at}\) and \(m_{1\nt}\) are not identified. Prove the bounds below.

\paragraph{The theorem to prove.}
Under monotonicity and double-index exclusion restriction, with bounded outcome
\(Y\in[\underline y,\overline y]\),
\[
\underline{\delta}\le \delta\le \overline{\delta},
\]
where
\[
\underline{\delta}
=
\delta'
-\overline y\,\pr(D=1\mid Z=0)
+\underline y\,\pr(D=0\mid Z=1),
\]
\[
\overline{\delta}
=
\delta'
-\underline y\,\pr(D=1\mid Z=0)
+\overline y\,\pr(D=0\mid Z=1),
\]
and
\[
\delta'=\E(DY\mid Z=1)-\E(Y-DY\mid Z=0).
\]

\paragraph{出处.}
Peng Ding, Chapter 22, Homework problems, ``Bounds on the average causal effect on the whole population''
\(\langle\)label: \ri{para::bounds-ACE}\(\rangle\).

\begin{prep}
  \item Identified latent-type proportions:
  \[
  \pi_{\at}=\pr(D=1\mid Z=0),\quad
  \pi_{\nt}=\pr(D=0\mid Z=1),\quad
  \pi_{\cp}=\E(D\mid Z=1)-\E(D\mid Z=0).
  \]
  \item Identified terms in \(\delta\):
  \[
  \pi_{\at}m_{1\at},\quad
  \pi_{\nt}m_{0\nt},\quad
  \pi_{\cp}m_{1\cp},\quad
  \pi_{\cp}m_{0\cp}.
  \]
  \item Unidentified terms are bounded by outcome support:
  \[
  m_{0\at}\in[\underline y,\overline y],
  \quad
  m_{1\nt}\in[\underline y,\overline y].
  \]
  \item Binary-outcome simplification:
  \[
  \underline{\delta}
  =
  \E(DY\mid Z=1)-\E(D+Y-DY\mid Z=0),
  \]
  \[
  \overline{\delta}
  =
  \E(DY+1-D\mid Z=1)-\E(Y-DY\mid Z=0).
  \]
\end{prep}

\begin{solutionbox}
\end{solutionbox}
\end{problemblock}

\begin{problemblock}{Chapter 23.1: More algebra for TSLS}
\paragraph{题面.}
\begin{enumerate}[label=(\alph*),leftmargin=2em]
  \item Show that \(\hat\Gamma\) in the first-stage OLS relation
  \(\hat D_i=\hat\Gamma\tran Z_i\) equals
  \[
  \hat{\Gamma}
  =
  \left(\sumn Z_iZ_i\tran\right)^{-1}\sumn Z_iD_i\tran.
  \]
  \item Show that \(\hat\beta_{\TSLS}\) reduces to \(\hat\beta_{\mathrm{iv}}\)
  if \(Z\) and \(D\) have the same dimension and
  \[
  n^{-1}\sumn Z_iZ_i\tran,
  \quad
  n^{-1}\sumn Z_iD_i\tran
  \]
  are both invertible.
\end{enumerate}

\paragraph{出处.}
Peng Ding, Chapter 23, Homework problems, ``More algebra for TSLS''
\(\langle\)label: \ri{prob::algebra-tsls}\(\rangle\).

\begin{prep}
  \item TSLS definition:
  first regress \(D\) on \(Z\) to get \(\hat D_i\), then regress \(Y\) on \(\hat D\).
  \item Explicit TSLS formula:
  \[
  \hat\beta_{\TSLS}
  =
  \left(\sumn \hat D_i\hat D_i\tran\right)^{-1}
  \sumn \hat D_iY_i.
  \]
  \item First-stage orthogonality:
  \[
  D_i=\hat D_i+\check D_i,
  \quad
  \sumn \hat D_i\check D_i\tran=0.
  \]
  \item Just-identified IV estimator:
  \[
  \hat\beta_{\mathrm{iv}}
  =
  \left(\sumn Z_iD_i\tran\right)^{-1}\sumn Z_iY_i.
  \]
  \item References: TSLS \citep{theil1953estimation,basmann1957generalized}.
\end{prep}

\begin{solutionbox}
\end{solutionbox}
\end{problemblock}

\begin{problemblock}{Chapter 23.2 optional: Equivalence between TSLS and ILS}
\paragraph{题面.}
Prove Theorem 23: with a single endogenous treatment and a single IV,
\[
\hat\beta_{1,\ILS}=\hat\beta_{1,\TSLS}.
\]
Remark in the book: use the FWL theorem.

\paragraph{出处.}
Peng Ding, Chapter 23, Homework problems, ``Equivalence between TSLS and ILS''
\(\langle\)label: \ri{para::tsls-ils}\(\rangle\). This problem is optional in the homework PDF.

\begin{prep}
  \item Structural equations:
  \[
  \begin{cases}
  Y_i=\beta_0+\beta_1D_i+\beta_2\tran X_i+\varepsilon_i,\\
  D_i=\gamma_0+\gamma_1Z_i+\gamma_2\tran X_i+\varepsilon_{2i}.
  \end{cases}
  \]
  \item Reduced form:
  \[
  \begin{cases}
  Y_i=\Gamma_0+\Gamma_1Z_i+\Gamma_2\tran X_i+\varepsilon_{1i},\\
  D_i=\gamma_0+\gamma_1Z_i+\gamma_2\tran X_i+\varepsilon_{2i}.
  \end{cases}
  \]
  \item ILS estimator:
  \[
  \hat\beta_{1,\ILS}=\frac{\hat\Gamma_1}{\hat\gamma_1}.
  \]
  \item The theorem is noted by \citet[][Section A.3]{imbens2014instrumental}.
\end{prep}

\begin{solutionbox}
\end{solutionbox}
\end{problemblock}

\begin{problemblock}{Chapter 23.3: Control function in the linear IV model}
\paragraph{题面.}
Define the control function estimator \(\hat\beta_{\textsc{cf}}\):
\begin{enumerate}[label=(\arabic*),leftmargin=2em]
  \item Run OLS of \(D\) on \(Z\), and obtain residuals \(\check D_i\).
  \item Run OLS of \(Y\) on \(D\) and \(\check D\), and take the coefficient of \(D\).
\end{enumerate}
Show that
\[
\hat\beta_{\textsc{cf}}=\hat\beta_{\TSLS}.
\]

\paragraph{出处.}
Peng Ding, Chapter 23, Homework problems, ``Control function in the linear instrumental variable model''
\(\langle\)label: \ri{para::control-function-iv}\(\rangle\).

\begin{prep}
  \item Decompose first stage:
  \[
  D=\hat D+\check D.
  \]
  \item OLS orthogonality implies \(\hat D\) is orthogonal to \(\check D\).
  \item Useful proof tools: FWL theorem, OLS orthogonality, and transformation of regressors.
  \item References: \citet{hausman1978specification} pointed out this result;
  \citet{wooldridge2015control} gives a broader discussion of control function methods.
\end{prep}

\begin{solutionbox}
\end{solutionbox}
\end{problemblock}

\section{Problem Set 5 额外题}

\begin{problemblock}{Problem 2: Wald and two-stage least squares}
\paragraph{题面.}
We derived the identification formula for CACE under randomization in class. Show that
\[
\CACE
=
\frac{\cov(Z_i,Y_i)}{\cov(Z_i,D_i)}.
\]

\begin{prep}
  \item Starting point from Chapter 21:
  \[
  \CACE
  =
  \frac{\E(Y\mid Z=1)-\E(Y\mid Z=0)}
  {\E(D\mid Z=1)-\E(D\mid Z=0)}.
  \]
  \item For binary \(Z\), with \(p=\pr(Z=1)\),
  \[
  \cov(Z,Y)=p(1-p)\{\E(Y\mid Z=1)-\E(Y\mid Z=0)\}.
  \]
  \[
  \cov(Z,D)=p(1-p)\{\E(D\mid Z=1)-\E(D\mid Z=0)\}.
  \]
  \item This also matches the single-IV linear model formula in Chapter 23:
  \[
  \beta=\frac{\cov(Z,Y)}{\cov(Z,D)}.
  \]
  \item References: Wald estimator \citep{wald1940fitting};
  IV/CACE framework \citep{imbens1994identification,angrist1996identification}.
\end{prep}

\begin{solutionbox}
\end{solutionbox}
\end{problemblock}

\begin{problemblock}{Problem 3 optional: Weighting method for CACE in observational studies}
\paragraph{题面.}
Replace randomization by conditional ignorability:
\[
Z_i\ind \{Y_i(1),Y_i(0),D_i(1),D_i(0)\}\mid X_i,
\]
while monotonicity and exclusion restriction continue to hold. Define
\[
\kappa_{0i}
=
(1-D_i)
\frac{(1-Z_i)-\pr(Z_i=0\mid X_i)}
{\pr(Z_i=0\mid X_i)\pr(Z_i=1\mid X_i)},
\]
\[
\kappa_{1i}
=
D_i
\frac{Z_i-\pr(Z_i=1\mid X_i)}
{\pr(Z_i=0\mid X_i)\pr(Z_i=1\mid X_i)}.
\]
Show that
\[
\E\{g(Y_i(0),X_i)\mid U_i=\cp\}
=
\frac{1}{\pr(U_i=\cp)}\E\{\kappa_{0i}g(Y,X)\},
\]
\[
\E\{g(Y_i(1),X_i)\mid U_i=\cp\}
=
\frac{1}{\pr(U_i=\cp)}\E\{\kappa_{1i}g(Y,X)\}.
\]

\begin{prep}
  \item Let \(e(X)=\pr(Z=1\mid X)\). Then
  \(\pr(Z=0\mid X)=1-e(X)\).
  \item Conditional ignorability allows replacing conditional potential-outcome expectations by observed conditional expectations within \(Z\)-groups.
  \item Monotonicity and exclusion restriction keep the same latent compliance interpretation as Chapter 21.
  \item This is the observational-study analogue of the mixture disentangling identities in Chapter 22.
  \item Useful target identity:
  \[
  \pr(U=\cp)=\E\{D(1)-D(0)\}.
  \]
\end{prep}

\begin{solutionbox}
\end{solutionbox}
\end{problemblock}

\begin{problemblock}{Problem 4: Analysis of return to schooling}
\paragraph{题面.}
\citet{card1993using} used geographic variation in college proximity to estimate the return to schooling.
Treatment is education, outcome is log wage, and the IV is college proximity.
The homework asks:
\begin{enumerate}[label=(\alph*),leftmargin=2em]
  \item Dichotomize education and infer the CACE for units with
  \(D(1)=1,D(0)=0\), first without covariates, then with covariates.
  Discuss threshold choice and sensitivity.
  \item Without dichotomizing education, use standard TSLS and compare results.
\end{enumerate}

\paragraph{出处.}
Homework PDF Problem 4; same data-analysis prompt is present as a commented problem in
Peng Ding Chapter 23 source. Reference: \citet{card1993using}.

\begin{prep}
  \item Causal quantity for dichotomized education:
  \[
  \tau_{\cp}(c)
  =
  \E\{Y(1)-Y(0)\mid D_c(1)=1,D_c(0)=0\},
  \]
  where \(D_c=1\{\text{education}\ge c\}\).
  \item Without covariates:
  \[
  \hat\tau_{\cp}(c)
  =
  \frac{\bar Y_{Z=1}-\bar Y_{Z=0}}
  {\bar D_{c,Z=1}-\bar D_{c,Z=0}}.
  \]
  \item With covariates: adjust both \(Y\) and dichotomized \(D_c\) for covariates,
  then take the ratio of the adjusted effects.
  \item TSLS with continuous education:
  \[
  \begin{cases}
  \text{education}_i=\gamma_0+\gamma_1Z_i+\gamma_2\tran X_i+\varepsilon_{2i},\\
  \log(\text{wage}_i)=\beta_0+\beta_1\text{education}_i+\beta_2\tran X_i+\varepsilon_i.
  \end{cases}
  \]
  \item Report assumptions: relevance \(\gamma_1\ne0\), IV exogeneity, exclusion restriction,
  monotonicity for the dichotomized CACE analysis, and model choice for TSLS.
  \item The current folder does not contain \(\ri{card.csv}\) or \(\ri{code\_bk.txt}\);
  add their local paths here before analysis:
  \[
  \ri{card.csv path: \underline{\hspace{6cm}}}
  \]
\end{prep}

\begin{solutionbox}
\end{solutionbox}
\end{problemblock}

\bibliographystyle{plainnat}
\bibliography{causal/causal}

\end{document}
